% ============================================================
%  conference_paper.tex
%  Automated Blood Group Detection Using Computer Vision and
%  Transfer Learning on Mobile-Optimized Deep Neural Networks
%
%  IEEE TENCON 2026 Conference Paper (4-5 pages)
%  Compile with:  pdflatex conference_paper.tex  (run twice)
% ============================================================
\documentclass[conference]{IEEEtran}

% ---- Packages -----------------------------------------------
\usepackage[T1]{fontenc}
\usepackage[utf8]{inputenc}
\usepackage{amsmath,amssymb}
\usepackage{graphicx}
\usepackage{booktabs}
\usepackage{array}
\usepackage{url}
\usepackage{cite}

% =============================================================
\begin{document}

% ---- Title & Author -----------------------------------------
\title{Automated Blood Group Detection Using Computer Vision and
       Transfer Learning on Mobile-Optimized Deep Neural Networks}

\author{\IEEEauthorblockN{Sayon Mitra}
\IEEEauthorblockA{Independent Researcher, AI and Biomedical Engineering\\
Email: sayon@sayonedu.in}}

\maketitle

% =============================================================
\begin{abstract}
Accurate and rapid blood group identification is critical in
transfusion medicine, emergency care, and surgical procedures.
Conventional haemagglutination-based typing is time-consuming,
operator-dependent, and unavailable in resource-limited settings.
This paper presents an automated, AI-driven blood group detection
system combining COCO bounding-box region extraction and MobileNetV2
transfer learning to classify ABO and Rh(D) blood groups from
photographs of standard blood-typing cards under ambient light.
The pipeline extracts three reagent zones (Anti-A, Anti-B, Anti-D),
classifies clumping intensity in each zone, and applies deterministic
ABO\,+\,Rh logic to produce the final blood group.
Trained on 6,456 images and evaluated on 383 test images, the system
achieves $\approx$94--95\,\% zone-level accuracy and a macro-average
F1-score of 0.947.
The lightweight MobileNetV2 backbone (14\,MB) supports deployment on
resource-constrained mobile devices, making the system suitable for
point-of-care screening in low-resource healthcare settings.
This work forms the AI core of \emph{Gyanbondhu}, an AI-powered
biodegradable blood group testing kit designed for field deployment.
This system is intended for assistive screening, not clinical
replacement.
\end{abstract}

\begin{IEEEkeywords}
blood group detection, transfer learning, MobileNetV2,
computer vision, point-of-care diagnostics, AI in healthcare
\end{IEEEkeywords}

% =============================================================
\section{Introduction}

Blood transfusion incompatibility remains a leading cause of
transfusion-related fatalities worldwide.
The ABO and Rh(D) systems are the most clinically significant
antigen systems; mismatching can trigger life-threatening haemolytic
reactions~\cite{lippi2011}.
Manual haemagglutination typing requires skilled personnel,
controlled conditions, and consistent lighting---resources unavailable
in rural clinics, military field hospitals, or emergency rooms.

The proliferation of high-resolution smartphone cameras and lightweight
deep learning frameworks creates an opportunity to democratise blood
group typing at the point of care.
This work presents a modular pipeline that classifies a blood group
from a single card photograph, substantially reducing the diagnostic
bottleneck in resource-limited settings.
It also forms the inference engine for \emph{Gyanbondhu}, an
AI-powered biodegradable blood group testing kit designed for
autonomous field deployment.

The key contributions of this paper are:
\begin{enumerate}
  \item An end-to-end pipeline: COCO region extraction
        $\rightarrow$ MobileNetV2 clumping classification
        $\rightarrow$ deterministic ABO\,+\,Rh decision logic.
  \item A curated annotated dataset of 6,456 training / 997
        validation / 383 test images achieving 94--95\,\%
        zone-level accuracy (macro-F1: 0.947).
  \item A mobile-deployable 14\,MB model with a TensorFlow Lite
        conversion roadmap for sub-50\,ms on-device inference.
\end{enumerate}

% =============================================================
\section{Related Work}

Haemagglutination-based typing has been the clinical gold standard
since Landsteiner's discovery of the ABO system~\cite{landsteiner1901},
but introduces inter-observer variability of 2--5\,\%~\cite{harmening2019}.
Early automated approaches used image heuristics: Saif~\emph{et al.}~\cite{saif2016}
achieved 88\,\% accuracy with HSV thresholding and blob detection;
Khan~\emph{et al.}~\cite{khan2018}
reported 91.3\,\% using an SVM on GLCM texture features.
CNN-based approaches raised the bar: ResNet-50~\cite{yildiz2020}
reached 93.7\,\% and VGG-16~\cite{ahmed2021} reached 94.2\,\%, but
both required pre-cropped single-reagent images under lab-controlled
conditions and did not address full-card region extraction.
Table~\ref{tab:comparison} shows that the proposed system is the
first to couple a complete card-level detection pipeline with a
mobile-friendly backbone trained on a dataset nearly five times
larger than prior work.

\begin{table}[htbp]
  \caption{Comparison with Prior Systems}
  \label{tab:comparison}
  \centering
  \renewcommand{\arraystretch}{1.1}
  \begin{tabular}{lp{1.5cm}cc}
    \toprule
    \textbf{System} & \textbf{Method} & \textbf{Acc.} & \textbf{Full Pipeline} \\
    \midrule
    Saif~\emph{et al.}~\cite{saif2016}    & HSV+Blob   & 88.0\,\% & No \\
    Khan~\emph{et al.}~\cite{khan2018}    & SVM+GLCM   & 91.3\,\% & No \\
    Yildiz~\emph{et al.}~\cite{yildiz2020}& ResNet-50  & 93.7\,\% & No \\
    Ahmed~\emph{et al.}~\cite{ahmed2021}  & VGG-16     & 94.2\,\% & No \\
    \textbf{Proposed}                     & \textbf{MNv2} & \textbf{94--95\,\%} & \textbf{Yes} \\
    \bottomrule
  \end{tabular}
\end{table}

% =============================================================
\section{Methodology}

\subsection{Dataset and Preprocessing}

Images of commercial disposable blood-typing cards were collected
and annotated using Roboflow in COCO JSON format~\cite{lin2014coco}.
Each full-card image contains three reagent spots: Anti-A, Anti-B,
and Anti-D.
The dataset covers all four ABO groups (A, B, AB, O) and both Rh
polarities, with statistics given in Table~\ref{tab:dataset}.

\begin{table}[htbp]
  \caption{Dataset Statistics}
  \label{tab:dataset}
  \centering
  \begin{tabular}{lcc}
    \toprule
    \textbf{Split} & \textbf{Images} & \textbf{Zone Crops (3$\times$)} \\
    \midrule
    Train      & 6,456 & 19,368 \\
    Validation & 997   & 2,991  \\
    Test       & 383   & 1,149  \\
    \midrule
    \textbf{Total} & \textbf{7,836} & \textbf{23,508} \\
    \bottomrule
  \end{tabular}
\end{table}

Training images are augmented with horizontal/vertical flip,
brightness ($\pm$0.15), contrast ($\pm$0.10), rotation
($\pm$15$^{\circ}$), and CLAHE normalisation to improve
generalisation across diverse lighting conditions and card brands.

\subsection{Region Extraction and Classification}

The COCO annotation file is parsed to retrieve bounding boxes for
the Anti-A, Anti-B, and Anti-D zones.
Each box is cropped and resized to $224\times224$ pixels, normalised
to $[0, 1]$, and labelled as one of three classes:
\textbf{Strong Clumping} (class 0), \textbf{Medium Clumping}
(class 1), or \textbf{No Clumping} (class 2).
Transfer learning is applied using MobileNetV2~\cite{sandler2018}
pretrained on ImageNet-1K, with the backbone frozen during the
initial phase and fine-tuned (last 30 layers, $\eta{=}10^{-5}$)
in a second phase.
The classification head appends GlobalAveragePooling2D (1280-d),
Dropout(0.3), Dense(128, ReLU), and Dense(3, Softmax).
Total trainable parameters in the head: $\approx$166\,K.
Training uses Adam ($\eta{=}10^{-4}$), sparse categorical
cross-entropy, batch size 32, early stopping (patience 3),
and learning-rate reduction (factor 0.5, patience 2).

\subsection{Blood Group Decision Logic}

Zone predictions are mapped to binary reactions ($R{=}1$ for clumping,
$R{=}0$ otherwise), with a confidence threshold of 0.90 applied to
the No-Clumping class to reduce false negatives from weak-D reactions:
\begin{equation}
  R =
  \begin{cases}
    1 & \text{class} \in \{0,1\}\\
    1 & \text{class} = 2 \;\text{and}\; P(\text{No Clump}) < 0.90 \\
    0 & \text{otherwise}
  \end{cases}
  \label{eq:reaction}
\end{equation}
The ABO group is derived from reactions $(R_A, R_B)$:
$A$ if $(+,-)$; $B$ if $(-,+)$; $AB$ if $(+,+)$; $O$ if $(-,-)$.
The Rh factor is $+$ if $R_D{=}1$, and $-$ if $R_D{=}0$.

% =============================================================
\section{Results}

\subsection{Zone-Level Performance}

Table~\ref{tab:perclass} gives per-class evaluation on 1,149 zone
crops from the 383 test images.
Strong Clumping achieves the highest F1 (0.957) due to its visually
distinct dense-aggregate texture.
Medium Clumping exhibits the lowest recall (0.914) because the
medium/strong boundary is physiologically continuous.
No Clumping is the primary source of Rh errors due to weak-D
reactions.

\begin{table}[htbp]
  \caption{Per-Class Evaluation (Zone-Level, Test Set)}
  \label{tab:perclass}
  \centering
  \renewcommand{\arraystretch}{1.1}
  \begin{tabular}{lcccc}
    \toprule
    \textbf{Class} & \textbf{Prec.} & \textbf{Rec.} &
    \textbf{F1} & \textbf{Support} \\
    \midrule
    Strong Clumping & 0.948 & 0.966 & 0.957 & 389 \\
    Medium Clumping & 0.925 & 0.914 & 0.919 & 357 \\
    No Clumping     & 0.968 & 0.961 & 0.964 & 403 \\
    \midrule
    \textbf{Macro Avg.} & \textbf{0.947} & \textbf{0.947} &
    \textbf{0.947} & \textbf{1,149} \\
    \bottomrule
  \end{tabular}
\end{table}

\subsection{End-to-End Accuracy}

End-to-end ABO+Rh accuracy on 383 test images is $\approx$91--93\,\%.
The 3\,pp shortfall relative to zone-level accuracy reflects error
compounding across three zones and ambiguity in weak-D (D\textsuperscript{u})
phenotype reactions.
Training and validation metrics are summarised in
Table~\ref{tab:results}.

\begin{table}[htbp]
  \caption{Training and Validation Performance}
  \label{tab:results}
  \centering
  \begin{tabular}{lc}
    \toprule
    \textbf{Metric} & \textbf{Value} \\
    \midrule
    Training Accuracy   & $\approx$97--98\,\% \\
    Validation Accuracy & $\approx$94--95\,\% \\
    Training Loss       & $\approx$0.07 \\
    Validation Loss     & $\approx$0.18 \\
    \bottomrule
  \end{tabular}
\end{table}

% =============================================================
\section{Discussion}

The proposed system outperforms all prior single-reagent approaches
(Table~\ref{tab:comparison}) while being the only method to handle
full-card images without manual region cropping.
Key strengths include: (i) a training corpus nearly five times
larger than prior work, exposing the model to diverse lighting
conditions and card brands; (ii) the frozen MobileNetV2 backbone,
which retains rich ImageNet representations and converges rapidly;
and (iii) a modular design that decouples region extraction,
classification, and decision logic.

Primary limitations are the current dependency on COCO bounding-box
annotations per image---requiring replacement by an automatic region
detector before unconstrained field deployment---and sensitivity to
extreme illumination.
Only MobileNetV2 was evaluated; comparative studies with
EfficientNet-B0 or ResNet-50 are identified as future work.
Uncertainty quantification (Monte Carlo Dropout or temperature
scaling) is absent; high-entropy predictions should trigger a
manual-review flag in safety-critical deployments.
\textbf{This system is intended for assistive screening, not
clinical replacement.}

% =============================================================
\section{Conclusion}

This paper presented a complete, annotation-driven blood group
detection pipeline combining COCO region extraction and MobileNetV2
transfer learning.
The system achieves $\approx$94--95\,\% zone-level accuracy and
$\approx$91--93\,\% end-to-end ABO+Rh accuracy (macro-F1: 0.947),
outperforming prior single-reagent systems with a fully automated
card-level pipeline.
The 14\,MB model supports deployment on mid-range smartphones via
TensorFlow Lite, making it viable for point-of-care screening in
low-resource settings.
This AI module forms the core of \emph{Gyanbondhu}, an AI-powered
biodegradable blood group testing kit enabling field-deployable
blood typing without laboratory infrastructure.
Future work will replace the COCO annotation dependency with an
on-device object detector (YOLOv8-nano), explore EfficientNet-B0
fine-tuning, and incorporate uncertainty quantification for
safety-critical deployment.

% =============================================================
\begin{thebibliography}{10}

\bibitem{lippi2011}
G.~Lippi, M.~Plebani, and E.~J.~Favaloro,
``Fatal errors in point-of-care testing,''
\emph{Clinica Chimica Acta}, vol.~412, no.~15--16,
pp.~1267--1273, 2011.

\bibitem{landsteiner1901}
K.~Landsteiner,
``Ueber Agglutinationserscheinungen normalen menschlichen Blutes,''
\emph{Wiener Klinische Wochenschrift}, vol.~14,
pp.~1132--1134, 1901.

\bibitem{harmening2019}
D.~M.~Harmening,
\emph{Modern Blood Banking and Transfusion Practices},
7th~ed. Philadelphia, PA, USA: F.\,A.~Davis Company, 2019.

\bibitem{saif2016}
A.~F.~M.~Saif \emph{et~al.},
``Automatic blood group identification using image processing,''
\emph{Int.\ J.\ Comput.\ Appl.}, vol.~140, no.~4,
pp.~1--5, 2016.

\bibitem{khan2018}
A.~Khan \emph{et~al.},
``Blood group detection using SVM with GLCM features,''
\emph{J.\ Med.\ Imaging Health Inf.}, vol.~8, no.~3,
pp.~614--620, 2018.

\bibitem{yildiz2020}
O.~Yildiz and M.~F.~Aslan,
``Blood type detection using deep convolutional neural networks,''
in \emph{Proc.\ IEEE INISTA 2020}, pp.~1--6.

\bibitem{ahmed2021}
S.~Ahmed \emph{et~al.},
``Deep learning approach for automated ABO/Rh blood group
determination,''
\emph{Appl.\ Sci.}, vol.~11, no.~14, p.~6660, 2021.

\bibitem{sandler2018}
M.~Sandler, A.~Howard, M.~Zhu, A.~Zhmoginov, and L.-C.~Chen,
``MobileNetV2: Inverted Residuals and Linear Bottlenecks,''
in \emph{Proc.\ IEEE/CVF CVPR 2018}, pp.~4510--4520.

\bibitem{lin2014coco}
T.-Y.~Lin \emph{et~al.},
``Microsoft COCO: Common Objects in Context,''
in \emph{Proc.\ ECCV 2014}, Lecture Notes Comput.\ Sci.,
vol.~8693. Springer, 2014, pp.~740--755.

\bibitem{howard2017mobilenets}
A.~G.~Howard \emph{et~al.},
``MobileNets: Efficient Convolutional Neural Networks for Mobile
Vision Applications,''
\emph{arXiv preprint arXiv:1704.04861}, 2017.

\end{thebibliography}

\end{document}
